\documentclass[12pt]{article}

% Set this switch to \showsolutiontrue to print the worked solutions.
\newif\ifshowsolution
\showsolutiontrue

% Update these fields before distributing the assignment.
\newcommand{\psnumber}{6}
\newcommand{\psdue}{TBD}
\usepackage[margin=1in]{geometry}
\usepackage[all]{xy}
\usepackage{enumitem}


\usepackage{amsmath,amsthm,amssymb,xcolor,latexsym}

\geometry{letterpaper}  
\usepackage{boxedminipage}
\usepackage{graphicx}
\usepackage{hyperref}
\hypersetup{
    colorlinks=true,
    citecolor=blue,
    linkcolor=blue,
    filecolor=cyan,
    urlcolor=magenta,
}
\usepackage[capitalize,nameinlink]{cleveref}

\newtheorem{problem}{Problem}
\newtheorem{lemma}{Lemma}


\usepackage{comment}
\ifshowsolution
  \newenvironment{solution}
    {\par\medskip\noindent\textbf{Solution.}\par\small}
    {\par}
\else
  \excludecomment{solution}
\fi
\newcommand{\pointsbox}[1]{%
  \colorbox{red!70}{\strut\textcolor{white}{\sffamily\bfseries\ \ #1\ \ }}%
}
\newcommand{\pagetag}[1]{%
  \colorbox{red!70}{\strut\textcolor{white}{\sffamily\bfseries\ \ #1\ \ }}%
}
\newcommand{\striptitle}[2]{%
  \noindent
  \begin{minipage}{0.07\linewidth}\centering \pagetag{#1}\end{minipage}%
  \hfill
  \begin{minipage}{0.89\linewidth}
    {\large\scshape #2}
  \end{minipage}\par\vspace{0.75em}
}
% --- Environments ---
\theoremstyle{plain}
\newtheorem{theorem}{Theorem}
\newtheorem{corollary}[theorem]{Corollary}
\theoremstyle{definition}
\newtheorem{definition}[theorem]{Definition}
\newtheorem{remark}[theorem]{Remark}

\newenvironment{boxedalgo}
  {\begin{center}\begin{boxedminipage}{1\linewidth}}
  {\end{boxedminipage}\end{center}}

\newenvironment{inlinealgo}[1]
  {\noindent\hrulefill\\\noindent{\bf {#1}}}
  {\hrulefill}


% --- Shortcuts ---
\newcommand{\E}{\mathbb{E}}
\newcommand{\Prob}{\mathbb{P}}
\newcommand{\R}{\mathbb{R}}
\newcommand{\inner}[2]{\left\langle #1,#2 \right\rangle}
\newcommand{\proj}{\mathrm{proj}}
\newcommand{\Id}{\mathbf{I}}
\newcommand{\ot}{\otimes}
\newcommand{\mc}{\mathcal}

\begin{document}
\noindent CS 578 Statistical Machine Learning (Fall 2026)\hfill Problem Set \# 6\\
\{\{Your name\}\} \hfill \textbf{Due: November 15th}

\noindent\hrulefill

\subsection*{Problem 1.}

Let $C_n$ ($n\ge3$) be an $n$-cycle with vertices $1,\dots,n$. Choose an independent set $S$ of $C_n$ uniformly at random, and let $q_n:=\Pr[1\in S]$. We showed in lecture (slide~35) that BP on $C_n$ converges to $m^{j\to i}_+=p$ on every directed edge, where $p:=(3-\sqrt5)/2=1/\varphi^2\approx0.382$ and $\varphi:=(1+\sqrt5)/2$. How does $p$ relate to $q_n$?

\begin{enumerate}[label=(\alph*),itemsep=0.3em]
\item Let $a_k$ be the number of independent sets of the $k$-vertex path, with $a_0:=1$. Derive a recurrence for $a_k$, and express $a_k$ in terms of the Fibonacci numbers ($F_0=0$, $F_1=1$, $F_{k+2}=F_{k+1}+F_k$).

\item By conditioning on whether vertex~$1$ is occupied, use (a) to find an exact formula for $q_n$. Evaluate $q_3$ and $q_4$.

\item At the BP fixed point, each neighbor $j\in\{2,n\}$ sends vertex~$1$ the message $m^{j\to1}_\pm=(p, 1-p)$. By the BP marginal formula (slide~34), the BP estimate of $q_n$ is
\[
\widehat q_{\mathrm{BP}}:=\frac{w_+}{w_++w_-},\qquad
w_\sigma:=\prod_{j\in\{2,n\}} \sum_{s\in\{\pm1\}}\psi_{1j}(\sigma,s)\,m^{j\to1}_s\quad(\sigma\in\{\pm1\}).
\]
Compute $w_+$, $w_-$, and $\widehat q_{\mathrm{BP}}$ as functions of $p$, and prove that $\lim_{n\to\infty}q_n=\widehat q_{\mathrm{BP}}$ (using $\lim_{k\to\infty}F_{k+1}/F_k=\varphi$).
\end{enumerate}
\noindent{\bfseries Solution.}\newline 
\   %%% <-- ERASE THIS LINE AND WRITE YOUR SOLUTION HERE
\newpage


\subsection*{Problem 2.}

Let $G=(V,E)$ be an unweighted, undirected graph, and let $E_v\subseteq E$ be the set of edges incident to $v\in V$. A set $M\subseteq E$ is a \emph{matching} if no two edges in $M$ share a node. We encode $M\subseteq E$ by $\mathbf{x}=(x_e)_{e\in E}\in\{-1,1\}^{E}$ with $x_e=+1$ if and only if $e\in M$, and write $M(\mathbf{x}):=\{e\in E:\, x_e=+1\}$.

\begin{enumerate}[label=(\alph*),itemsep=0.3em]
\item Let $\mu$ be the uniform distribution over all matchings of $G$. Express $\mu$ as a Gibbs measure of the form
\[
\mu(\mathbf{x})=\frac{1}{Z}\prod_{(e,f)\in F}\psi_{ef}(x_e,x_f),\qquad \mathbf{x}\in\{-1,1\}^{E},
\]
by specifying the interaction set $F$ and the compatibility functions $\psi_{ef}$.

\item For $\mu$ as in (a), write the BP fixed-point equations (slide~34) as a closed system of equations for the unknowns $m^{e\to f}_+$. Then compute the Bethe free energy $G_\beta$ at a BP fixed point (slide~43) \emph{explicitly}: evaluate $Z_e$, $Z_{ef}$, and $Z_{e;f}$ in terms of $m^{e\to f}_+$, substitute them into $G_\beta$, and simplify so that your final formula contains only the $m^{e\to f}_+$.

\item Now consider the weighted matching distribution $\mu_\lambda(\mathbf{x})\propto\lambda^{|M(\mathbf{x})|}\,\mathbf{1}\{M(\mathbf{x})\text{ is a matching}\}$ with $\lambda>1$, which assigns larger weight to larger matchings. Assuming $G$ is a connected graph and has at least two edges, repeat (a) and (b) for $\mu_\lambda$. 
\end{enumerate}

\noindent{\bfseries Solution.}\newline 
\   %%% <-- ERASE THIS LINE AND WRITE YOUR SOLUTION HERE
\newpage
\end{document}